\section{未锚定账本事件的叙事补充：southern-song}

\label{sec:anchor-batch-two-southern-song}

本组为尚无正文跳转目标的账本事件补入锚点。每一记录置于战争、行政、环境、迁徙与地方社会的关系中，并按来源限度约束叙事。

\subsection{制度、区域与材料边界}

\eventanchor{SS-1194-EVENT-129}{1194年·宁宗即位，韩侂胄等在中枢取得影响}账本条目记录南宋在1194年的“宁宗即位，韩侂胄等在中枢取得影响”。其背景是前期边防、财政和统治联盟的压力构成直接背景，具体条件须按本条材料分辨，后续影响为改变后续资源配置与政治选择；实际执行范围和社会影响仍须分项核验。政权、军队或制度的变化须经由道路、文书、资源和地方中介才会落实，城市、乡村、边地和流动人群不可能在同一时间承受相同结果。

本条使用《宋史》卷24—47〈本纪〉；《建炎以来系年要录》本年条；《宋会要辑稿》相关门类，并标示“编年主干可靠；细节、规模与动机须与对方材料、地方文书或考古材料互校”。这些材料能够钉住年序、命令、战役或局部地点，却不能直接推出人口、税收、身份和社会动机的总量。应区分诏令与实施、条约与边境生活、军报与伤亡统计，也不将官府的族群或行政称谓写成固定的社会本质。

从区域史角度看，该事件是资源、信息与权力重新连接的节点。不同地方会因市场、交通、宗教组织、家庭策略和环境条件而产生不同反应；材料的沉默限制我们对未被记录者所能作出的判断。

\eventanchor{SS-1194-REGNAL-YEAR}{1194年·南宋以宋光宗在位第5年作为诏令、赋役与外交文书的纪年框架}账本条目记录南宋在1194年的“南宋以宋光宗在位第5年作为诏令、赋役与外交文书的纪年框架”。其背景是新岁颁历、年号和纪年是中枢与地方行政沟通的共同前提，后续影响为为同年政令、战事和地方材料提供日期锚点，不能单独推知具体政策执行。政权、军队或制度的变化须经由道路、文书、资源和地方中介才会落实，城市、乡村、边地和流动人群不可能在同一时间承受相同结果。

本条使用《宋史》卷24—47〈本纪〉；《建炎以来系年要录》本年条；《宋会要辑稿》相关门类，并标示“纪年框架可靠；该记录仅确证政权纪年连续，年内具体事项须由本年条另证”。这些材料能够钉住年序、命令、战役或局部地点，却不能直接推出人口、税收、身份和社会动机的总量。应区分诏令与实施、条约与边境生活、军报与伤亡统计，也不将官府的族群或行政称谓写成固定的社会本质。

从区域史角度看，该事件是资源、信息与权力重新连接的节点。不同地方会因市场、交通、宗教组织、家庭策略和环境条件而产生不同反应；材料的沉默限制我们对未被记录者所能作出的判断。

\eventanchor{SS-1195-REGNAL-YEAR}{1195年·南宋以宋宁宗在位第1年作为诏令、赋役与外交文书的纪年框架}账本条目记录南宋在1195年的“南宋以宋宁宗在位第1年作为诏令、赋役与外交文书的纪年框架”。其背景是新岁颁历、年号和纪年是中枢与地方行政沟通的共同前提，后续影响为为同年政令、战事和地方材料提供日期锚点，不能单独推知具体政策执行。政权、军队或制度的变化须经由道路、文书、资源和地方中介才会落实，城市、乡村、边地和流动人群不可能在同一时间承受相同结果。

本条使用《宋史》卷24—47〈本纪〉；《建炎以来系年要录》本年条；《宋会要辑稿》相关门类，并标示“纪年框架可靠；该记录仅确证政权纪年连续，年内具体事项须由本年条另证”。这些材料能够钉住年序、命令、战役或局部地点，却不能直接推出人口、税收、身份和社会动机的总量。应区分诏令与实施、条约与边境生活、军报与伤亡统计，也不将官府的族群或行政称谓写成固定的社会本质。

从区域史角度看，该事件是资源、信息与权力重新连接的节点。不同地方会因市场、交通、宗教组织、家庭策略和环境条件而产生不同反应；材料的沉默限制我们对未被记录者所能作出的判断。

\eventanchor{SS-1196-EVENT-130}{1196年·庆元党禁限制朱熹及相关士人的政治与学术活动}账本条目记录南宋在1196年的“庆元党禁限制朱熹及相关士人的政治与学术活动”。其背景是前期边防、财政和统治联盟的压力构成直接背景，具体条件须按本条材料分辨，后续影响为改变后续资源配置与政治选择；实际执行范围和社会影响仍须分项核验。政权、军队或制度的变化须经由道路、文书、资源和地方中介才会落实，城市、乡村、边地和流动人群不可能在同一时间承受相同结果。

本条使用《宋史》卷24—47〈本纪〉；《建炎以来系年要录》本年条；《宋会要辑稿》相关门类，并标示“编年主干可靠；细节、规模与动机须与对方材料、地方文书或考古材料互校”。这些材料能够钉住年序、命令、战役或局部地点，却不能直接推出人口、税收、身份和社会动机的总量。应区分诏令与实施、条约与边境生活、军报与伤亡统计，也不将官府的族群或行政称谓写成固定的社会本质。

从区域史角度看，该事件是资源、信息与权力重新连接的节点。不同地方会因市场、交通、宗教组织、家庭策略和环境条件而产生不同反应；材料的沉默限制我们对未被记录者所能作出的判断。

\eventanchor{SS-1196-REGNAL-YEAR}{1196年·南宋以宋宁宗在位第2年作为诏令、赋役与外交文书的纪年框架}账本条目记录南宋在1196年的“南宋以宋宁宗在位第2年作为诏令、赋役与外交文书的纪年框架”。其背景是新岁颁历、年号和纪年是中枢与地方行政沟通的共同前提，后续影响为为同年政令、战事和地方材料提供日期锚点，不能单独推知具体政策执行。政权、军队或制度的变化须经由道路、文书、资源和地方中介才会落实，城市、乡村、边地和流动人群不可能在同一时间承受相同结果。

本条使用《宋史》卷24—47〈本纪〉；《建炎以来系年要录》本年条；《宋会要辑稿》相关门类，并标示“纪年框架可靠；该记录仅确证政权纪年连续，年内具体事项须由本年条另证”。这些材料能够钉住年序、命令、战役或局部地点，却不能直接推出人口、税收、身份和社会动机的总量。应区分诏令与实施、条约与边境生活、军报与伤亡统计，也不将官府的族群或行政称谓写成固定的社会本质。

从区域史角度看，该事件是资源、信息与权力重新连接的节点。不同地方会因市场、交通、宗教组织、家庭策略和环境条件而产生不同反应；材料的沉默限制我们对未被记录者所能作出的判断。

\eventanchor{SS-1197-REGNAL-YEAR}{1197年·南宋以宋宁宗在位第3年作为诏令、赋役与外交文书的纪年框架}账本条目记录南宋在1197年的“南宋以宋宁宗在位第3年作为诏令、赋役与外交文书的纪年框架”。其背景是新岁颁历、年号和纪年是中枢与地方行政沟通的共同前提，后续影响为为同年政令、战事和地方材料提供日期锚点，不能单独推知具体政策执行。政权、军队或制度的变化须经由道路、文书、资源和地方中介才会落实，城市、乡村、边地和流动人群不可能在同一时间承受相同结果。

本条使用《宋史》卷24—47〈本纪〉；《建炎以来系年要录》本年条；《宋会要辑稿》相关门类，并标示“纪年框架可靠；该记录仅确证政权纪年连续，年内具体事项须由本年条另证”。这些材料能够钉住年序、命令、战役或局部地点，却不能直接推出人口、税收、身份和社会动机的总量。应区分诏令与实施、条约与边境生活、军报与伤亡统计，也不将官府的族群或行政称谓写成固定的社会本质。

从区域史角度看，该事件是资源、信息与权力重新连接的节点。不同地方会因市场、交通、宗教组织、家庭策略和环境条件而产生不同反应；材料的沉默限制我们对未被记录者所能作出的判断。

\eventanchor{SS-1198-REGNAL-YEAR}{1198年·南宋以宋宁宗在位第4年作为诏令、赋役与外交文书的纪年框架}账本条目记录南宋在1198年的“南宋以宋宁宗在位第4年作为诏令、赋役与外交文书的纪年框架”。其背景是新岁颁历、年号和纪年是中枢与地方行政沟通的共同前提，后续影响为为同年政令、战事和地方材料提供日期锚点，不能单独推知具体政策执行。政权、军队或制度的变化须经由道路、文书、资源和地方中介才会落实，城市、乡村、边地和流动人群不可能在同一时间承受相同结果。

本条使用《宋史》卷24—47〈本纪〉；《建炎以来系年要录》本年条；《宋会要辑稿》相关门类，并标示“纪年框架可靠；该记录仅确证政权纪年连续，年内具体事项须由本年条另证”。这些材料能够钉住年序、命令、战役或局部地点，却不能直接推出人口、税收、身份和社会动机的总量。应区分诏令与实施、条约与边境生活、军报与伤亡统计，也不将官府的族群或行政称谓写成固定的社会本质。

从区域史角度看，该事件是资源、信息与权力重新连接的节点。不同地方会因市场、交通、宗教组织、家庭策略和环境条件而产生不同反应；材料的沉默限制我们对未被记录者所能作出的判断。

\eventanchor{SS-1199-REGNAL-YEAR}{1199年·南宋以宋宁宗在位第5年作为诏令、赋役与外交文书的纪年框架}账本条目记录南宋在1199年的“南宋以宋宁宗在位第5年作为诏令、赋役与外交文书的纪年框架”。其背景是新岁颁历、年号和纪年是中枢与地方行政沟通的共同前提，后续影响为为同年政令、战事和地方材料提供日期锚点，不能单独推知具体政策执行。政权、军队或制度的变化须经由道路、文书、资源和地方中介才会落实，城市、乡村、边地和流动人群不可能在同一时间承受相同结果。

本条使用《宋史》卷24—47〈本纪〉；《建炎以来系年要录》本年条；《宋会要辑稿》相关门类，并标示“纪年框架可靠；该记录仅确证政权纪年连续，年内具体事项须由本年条另证”。这些材料能够钉住年序、命令、战役或局部地点，却不能直接推出人口、税收、身份和社会动机的总量。应区分诏令与实施、条约与边境生活、军报与伤亡统计，也不将官府的族群或行政称谓写成固定的社会本质。

从区域史角度看，该事件是资源、信息与权力重新连接的节点。不同地方会因市场、交通、宗教组织、家庭策略和环境条件而产生不同反应；材料的沉默限制我们对未被记录者所能作出的判断。

\eventanchor{SS-1200-REGNAL-YEAR}{1200年·南宋以宋宁宗在位第6年作为诏令、赋役与外交文书的纪年框架}账本条目记录南宋在1200年的“南宋以宋宁宗在位第6年作为诏令、赋役与外交文书的纪年框架”。其背景是新岁颁历、年号和纪年是中枢与地方行政沟通的共同前提，后续影响为为同年政令、战事和地方材料提供日期锚点，不能单独推知具体政策执行。政权、军队或制度的变化须经由道路、文书、资源和地方中介才会落实，城市、乡村、边地和流动人群不可能在同一时间承受相同结果。

本条使用《宋史》卷24—47〈本纪〉；《建炎以来系年要录》本年条；《宋会要辑稿》相关门类，并标示“纪年框架可靠；该记录仅确证政权纪年连续，年内具体事项须由本年条另证”。这些材料能够钉住年序、命令、战役或局部地点，却不能直接推出人口、税收、身份和社会动机的总量。应区分诏令与实施、条约与边境生活、军报与伤亡统计，也不将官府的族群或行政称谓写成固定的社会本质。

从区域史角度看，该事件是资源、信息与权力重新连接的节点。不同地方会因市场、交通、宗教组织、家庭策略和环境条件而产生不同反应；材料的沉默限制我们对未被记录者所能作出的判断。

\eventanchor{SS-1201-REGNAL-YEAR}{1201年·南宋以宋宁宗在位第7年作为诏令、赋役与外交文书的纪年框架}账本条目记录南宋在1201年的“南宋以宋宁宗在位第7年作为诏令、赋役与外交文书的纪年框架”。其背景是新岁颁历、年号和纪年是中枢与地方行政沟通的共同前提，后续影响为为同年政令、战事和地方材料提供日期锚点，不能单独推知具体政策执行。政权、军队或制度的变化须经由道路、文书、资源和地方中介才会落实，城市、乡村、边地和流动人群不可能在同一时间承受相同结果。

本条使用《宋史》卷24—47〈本纪〉；《建炎以来系年要录》本年条；《宋会要辑稿》相关门类，并标示“纪年框架可靠；该记录仅确证政权纪年连续，年内具体事项须由本年条另证”。这些材料能够钉住年序、命令、战役或局部地点，却不能直接推出人口、税收、身份和社会动机的总量。应区分诏令与实施、条约与边境生活、军报与伤亡统计，也不将官府的族群或行政称谓写成固定的社会本质。

从区域史角度看，该事件是资源、信息与权力重新连接的节点。不同地方会因市场、交通、宗教组织、家庭策略和环境条件而产生不同反应；材料的沉默限制我们对未被记录者所能作出的判断。

\eventanchor{SS-1202-REGNAL-YEAR}{1202年·南宋以宋宁宗在位第8年作为诏令、赋役与外交文书的纪年框架}账本条目记录南宋在1202年的“南宋以宋宁宗在位第8年作为诏令、赋役与外交文书的纪年框架”。其背景是新岁颁历、年号和纪年是中枢与地方行政沟通的共同前提，后续影响为为同年政令、战事和地方材料提供日期锚点，不能单独推知具体政策执行。政权、军队或制度的变化须经由道路、文书、资源和地方中介才会落实，城市、乡村、边地和流动人群不可能在同一时间承受相同结果。

本条使用《宋史》卷24—47〈本纪〉；《建炎以来系年要录》本年条；《宋会要辑稿》相关门类，并标示“纪年框架可靠；该记录仅确证政权纪年连续，年内具体事项须由本年条另证”。这些材料能够钉住年序、命令、战役或局部地点，却不能直接推出人口、税收、身份和社会动机的总量。应区分诏令与实施、条约与边境生活、军报与伤亡统计，也不将官府的族群或行政称谓写成固定的社会本质。

从区域史角度看，该事件是资源、信息与权力重新连接的节点。不同地方会因市场、交通、宗教组织、家庭策略和环境条件而产生不同反应；材料的沉默限制我们对未被记录者所能作出的判断。

\eventanchor{SS-1203-REGNAL-YEAR}{1203年·南宋以宋宁宗在位第9年作为诏令、赋役与外交文书的纪年框架}账本条目记录南宋在1203年的“南宋以宋宁宗在位第9年作为诏令、赋役与外交文书的纪年框架”。其背景是新岁颁历、年号和纪年是中枢与地方行政沟通的共同前提，后续影响为为同年政令、战事和地方材料提供日期锚点，不能单独推知具体政策执行。政权、军队或制度的变化须经由道路、文书、资源和地方中介才会落实，城市、乡村、边地和流动人群不可能在同一时间承受相同结果。

本条使用《宋史》卷24—47〈本纪〉；《建炎以来系年要录》本年条；《宋会要辑稿》相关门类，并标示“纪年框架可靠；该记录仅确证政权纪年连续，年内具体事项须由本年条另证”。这些材料能够钉住年序、命令、战役或局部地点，却不能直接推出人口、税收、身份和社会动机的总量。应区分诏令与实施、条约与边境生活、军报与伤亡统计，也不将官府的族群或行政称谓写成固定的社会本质。

从区域史角度看，该事件是资源、信息与权力重新连接的节点。不同地方会因市场、交通、宗教组织、家庭策略和环境条件而产生不同反应；材料的沉默限制我们对未被记录者所能作出的判断。

\eventanchor{SS-1204-EVENT-131}{1204年·韩侂胄推动北伐准备}账本条目记录南宋与金在1204年的“韩侂胄推动北伐准备”。其背景是前期边防、财政和统治联盟的压力构成直接背景，具体条件须按本条材料分辨，后续影响为改变后续资源配置与政治选择；实际执行范围和社会影响仍须分项核验。政权、军队或制度的变化须经由道路、文书、资源和地方中介才会落实，城市、乡村、边地和流动人群不可能在同一时间承受相同结果。

本条使用《宋史》卷24—47〈本纪〉；《建炎以来系年要录》本年条；《宋会要辑稿》相关门类，并标示“编年主干可靠；细节、规模与动机须与对方材料、地方文书或考古材料互校”。这些材料能够钉住年序、命令、战役或局部地点，却不能直接推出人口、税收、身份和社会动机的总量。应区分诏令与实施、条约与边境生活、军报与伤亡统计，也不将官府的族群或行政称谓写成固定的社会本质。

从区域史角度看，该事件是资源、信息与权力重新连接的节点。不同地方会因市场、交通、宗教组织、家庭策略和环境条件而产生不同反应；材料的沉默限制我们对未被记录者所能作出的判断。

\eventanchor{SS-1204-REGNAL-YEAR}{1204年·南宋以宋宁宗在位第10年作为诏令、赋役与外交文书的纪年框架}账本条目记录南宋在1204年的“南宋以宋宁宗在位第10年作为诏令、赋役与外交文书的纪年框架”。其背景是新岁颁历、年号和纪年是中枢与地方行政沟通的共同前提，后续影响为为同年政令、战事和地方材料提供日期锚点，不能单独推知具体政策执行。政权、军队或制度的变化须经由道路、文书、资源和地方中介才会落实，城市、乡村、边地和流动人群不可能在同一时间承受相同结果。

本条使用《宋史》卷24—47〈本纪〉；《建炎以来系年要录》本年条；《宋会要辑稿》相关门类，并标示“纪年框架可靠；该记录仅确证政权纪年连续，年内具体事项须由本年条另证”。这些材料能够钉住年序、命令、战役或局部地点，却不能直接推出人口、税收、身份和社会动机的总量。应区分诏令与实施、条约与边境生活、军报与伤亡统计，也不将官府的族群或行政称谓写成固定的社会本质。

从区域史角度看，该事件是资源、信息与权力重新连接的节点。不同地方会因市场、交通、宗教组织、家庭策略和环境条件而产生不同反应；材料的沉默限制我们对未被记录者所能作出的判断。

\eventanchor{SS-1205-REGNAL-YEAR}{1205年·南宋以宋宁宗在位第11年作为诏令、赋役与外交文书的纪年框架}账本条目记录南宋在1205年的“南宋以宋宁宗在位第11年作为诏令、赋役与外交文书的纪年框架”。其背景是新岁颁历、年号和纪年是中枢与地方行政沟通的共同前提，后续影响为为同年政令、战事和地方材料提供日期锚点，不能单独推知具体政策执行。政权、军队或制度的变化须经由道路、文书、资源和地方中介才会落实，城市、乡村、边地和流动人群不可能在同一时间承受相同结果。

本条使用《宋史》卷24—47〈本纪〉；《建炎以来系年要录》本年条；《宋会要辑稿》相关门类，并标示“纪年框架可靠；该记录仅确证政权纪年连续，年内具体事项须由本年条另证”。这些材料能够钉住年序、命令、战役或局部地点，却不能直接推出人口、税收、身份和社会动机的总量。应区分诏令与实施、条约与边境生活、军报与伤亡统计，也不将官府的族群或行政称谓写成固定的社会本质。

从区域史角度看，该事件是资源、信息与权力重新连接的节点。不同地方会因市场、交通、宗教组织、家庭策略和环境条件而产生不同反应；材料的沉默限制我们对未被记录者所能作出的判断。

\eventanchor{SS-1206-EVENT-132}{1206年·南宋发动开禧北伐，战事不利}账本条目记录南宋与金在1206年的“南宋发动开禧北伐，战事不利”。其背景是前期边防、财政和统治联盟的压力构成直接背景，具体条件须按本条材料分辨，后续影响为改变后续资源配置与政治选择；实际执行范围和社会影响仍须分项核验。政权、军队或制度的变化须经由道路、文书、资源和地方中介才会落实，城市、乡村、边地和流动人群不可能在同一时间承受相同结果。

本条使用《宋史》卷24—47〈本纪〉；《建炎以来系年要录》本年条；《宋会要辑稿》相关门类，并标示“编年主干可靠；细节、规模与动机须与对方材料、地方文书或考古材料互校”。这些材料能够钉住年序、命令、战役或局部地点，却不能直接推出人口、税收、身份和社会动机的总量。应区分诏令与实施、条约与边境生活、军报与伤亡统计，也不将官府的族群或行政称谓写成固定的社会本质。

从区域史角度看，该事件是资源、信息与权力重新连接的节点。不同地方会因市场、交通、宗教组织、家庭策略和环境条件而产生不同反应；材料的沉默限制我们对未被记录者所能作出的判断。

\eventanchor{SS-1206-REGNAL-YEAR}{1206年·南宋以宋宁宗在位第12年作为诏令、赋役与外交文书的纪年框架}账本条目记录南宋在1206年的“南宋以宋宁宗在位第12年作为诏令、赋役与外交文书的纪年框架”。其背景是新岁颁历、年号和纪年是中枢与地方行政沟通的共同前提，后续影响为为同年政令、战事和地方材料提供日期锚点，不能单独推知具体政策执行。政权、军队或制度的变化须经由道路、文书、资源和地方中介才会落实，城市、乡村、边地和流动人群不可能在同一时间承受相同结果。

本条使用《宋史》卷24—47〈本纪〉；《建炎以来系年要录》本年条；《宋会要辑稿》相关门类，并标示“纪年框架可靠；该记录仅确证政权纪年连续，年内具体事项须由本年条另证”。这些材料能够钉住年序、命令、战役或局部地点，却不能直接推出人口、税收、身份和社会动机的总量。应区分诏令与实施、条约与边境生活、军报与伤亡统计，也不将官府的族群或行政称谓写成固定的社会本质。

从区域史角度看，该事件是资源、信息与权力重新连接的节点。不同地方会因市场、交通、宗教组织、家庭策略和环境条件而产生不同反应；材料的沉默限制我们对未被记录者所能作出的判断。

\eventanchor{SS-1207-EVENT-133}{1207年·韩侂胄被杀，朝廷转向结束北伐}账本条目记录南宋在1207年的“韩侂胄被杀，朝廷转向结束北伐”。其背景是前期边防、财政和统治联盟的压力构成直接背景，具体条件须按本条材料分辨，后续影响为改变后续资源配置与政治选择；实际执行范围和社会影响仍须分项核验。政权、军队或制度的变化须经由道路、文书、资源和地方中介才会落实，城市、乡村、边地和流动人群不可能在同一时间承受相同结果。

本条使用《宋史》卷24—47〈本纪〉；《建炎以来系年要录》本年条；《宋会要辑稿》相关门类，并标示“编年主干可靠；细节、规模与动机须与对方材料、地方文书或考古材料互校”。这些材料能够钉住年序、命令、战役或局部地点，却不能直接推出人口、税收、身份和社会动机的总量。应区分诏令与实施、条约与边境生活、军报与伤亡统计，也不将官府的族群或行政称谓写成固定的社会本质。

从区域史角度看，该事件是资源、信息与权力重新连接的节点。不同地方会因市场、交通、宗教组织、家庭策略和环境条件而产生不同反应；材料的沉默限制我们对未被记录者所能作出的判断。

\eventanchor{SS-1207-REGNAL-YEAR}{1207年·南宋以宋宁宗在位第13年作为诏令、赋役与外交文书的纪年框架}账本条目记录南宋在1207年的“南宋以宋宁宗在位第13年作为诏令、赋役与外交文书的纪年框架”。其背景是新岁颁历、年号和纪年是中枢与地方行政沟通的共同前提，后续影响为为同年政令、战事和地方材料提供日期锚点，不能单独推知具体政策执行。政权、军队或制度的变化须经由道路、文书、资源和地方中介才会落实，城市、乡村、边地和流动人群不可能在同一时间承受相同结果。

本条使用《宋史》卷24—47〈本纪〉；《建炎以来系年要录》本年条；《宋会要辑稿》相关门类，并标示“纪年框架可靠；该记录仅确证政权纪年连续，年内具体事项须由本年条另证”。这些材料能够钉住年序、命令、战役或局部地点，却不能直接推出人口、税收、身份和社会动机的总量。应区分诏令与实施、条约与边境生活、军报与伤亡统计，也不将官府的族群或行政称谓写成固定的社会本质。

从区域史角度看，该事件是资源、信息与权力重新连接的节点。不同地方会因市场、交通、宗教组织、家庭策略和环境条件而产生不同反应；材料的沉默限制我们对未被记录者所能作出的判断。

\eventanchor{SS-1208-EVENT-134}{1208年·南宋与金订立嘉定和议}账本条目记录南宋与金在1208年的“南宋与金订立嘉定和议”。其背景是前期边防、财政和统治联盟的压力构成直接背景，具体条件须按本条材料分辨，后续影响为改变后续资源配置与政治选择；实际执行范围和社会影响仍须分项核验。政权、军队或制度的变化须经由道路、文书、资源和地方中介才会落实，城市、乡村、边地和流动人群不可能在同一时间承受相同结果。

本条使用《宋史》卷24—47〈本纪〉；《建炎以来系年要录》本年条；《宋会要辑稿》相关门类，并标示“编年主干可靠；细节、规模与动机须与对方材料、地方文书或考古材料互校”。这些材料能够钉住年序、命令、战役或局部地点，却不能直接推出人口、税收、身份和社会动机的总量。应区分诏令与实施、条约与边境生活、军报与伤亡统计，也不将官府的族群或行政称谓写成固定的社会本质。

从区域史角度看，该事件是资源、信息与权力重新连接的节点。不同地方会因市场、交通、宗教组织、家庭策略和环境条件而产生不同反应；材料的沉默限制我们对未被记录者所能作出的判断。

\eventanchor{SS-1208-REGNAL-YEAR}{1208年·南宋以宋宁宗在位第14年作为诏令、赋役与外交文书的纪年框架}账本条目记录南宋在1208年的“南宋以宋宁宗在位第14年作为诏令、赋役与外交文书的纪年框架”。其背景是新岁颁历、年号和纪年是中枢与地方行政沟通的共同前提，后续影响为为同年政令、战事和地方材料提供日期锚点，不能单独推知具体政策执行。政权、军队或制度的变化须经由道路、文书、资源和地方中介才会落实，城市、乡村、边地和流动人群不可能在同一时间承受相同结果。

本条使用《宋史》卷24—47〈本纪〉；《建炎以来系年要录》本年条；《宋会要辑稿》相关门类，并标示“纪年框架可靠；该记录仅确证政权纪年连续，年内具体事项须由本年条另证”。这些材料能够钉住年序、命令、战役或局部地点，却不能直接推出人口、税收、身份和社会动机的总量。应区分诏令与实施、条约与边境生活、军报与伤亡统计，也不将官府的族群或行政称谓写成固定的社会本质。

从区域史角度看，该事件是资源、信息与权力重新连接的节点。不同地方会因市场、交通、宗教组织、家庭策略和环境条件而产生不同反应；材料的沉默限制我们对未被记录者所能作出的判断。

\eventanchor{SS-1209-REGNAL-YEAR}{1209年·南宋以宋宁宗在位第15年作为诏令、赋役与外交文书的纪年框架}账本条目记录南宋在1209年的“南宋以宋宁宗在位第15年作为诏令、赋役与外交文书的纪年框架”。其背景是新岁颁历、年号和纪年是中枢与地方行政沟通的共同前提，后续影响为为同年政令、战事和地方材料提供日期锚点，不能单独推知具体政策执行。政权、军队或制度的变化须经由道路、文书、资源和地方中介才会落实，城市、乡村、边地和流动人群不可能在同一时间承受相同结果。

本条使用《宋史》卷24—47〈本纪〉；《建炎以来系年要录》本年条；《宋会要辑稿》相关门类，并标示“纪年框架可靠；该记录仅确证政权纪年连续，年内具体事项须由本年条另证”。这些材料能够钉住年序、命令、战役或局部地点，却不能直接推出人口、税收、身份和社会动机的总量。应区分诏令与实施、条约与边境生活、军报与伤亡统计，也不将官府的族群或行政称谓写成固定的社会本质。

从区域史角度看，该事件是资源、信息与权力重新连接的节点。不同地方会因市场、交通、宗教组织、家庭策略和环境条件而产生不同反应；材料的沉默限制我们对未被记录者所能作出的判断。

\eventanchor{SS-1210-REGNAL-YEAR}{1210年·南宋以宋宁宗在位第16年作为诏令、赋役与外交文书的纪年框架}账本条目记录南宋在1210年的“南宋以宋宁宗在位第16年作为诏令、赋役与外交文书的纪年框架”。其背景是新岁颁历、年号和纪年是中枢与地方行政沟通的共同前提，后续影响为为同年政令、战事和地方材料提供日期锚点，不能单独推知具体政策执行。政权、军队或制度的变化须经由道路、文书、资源和地方中介才会落实，城市、乡村、边地和流动人群不可能在同一时间承受相同结果。

本条使用《宋史》卷24—47〈本纪〉；《建炎以来系年要录》本年条；《宋会要辑稿》相关门类，并标示“纪年框架可靠；该记录仅确证政权纪年连续，年内具体事项须由本年条另证”。这些材料能够钉住年序、命令、战役或局部地点，却不能直接推出人口、税收、身份和社会动机的总量。应区分诏令与实施、条约与边境生活、军报与伤亡统计，也不将官府的族群或行政称谓写成固定的社会本质。

从区域史角度看，该事件是资源、信息与权力重新连接的节点。不同地方会因市场、交通、宗教组织、家庭策略和环境条件而产生不同反应；材料的沉默限制我们对未被记录者所能作出的判断。

\eventanchor{SS-1211-REGNAL-YEAR}{1211年·南宋以宋宁宗在位第17年作为诏令、赋役与外交文书的纪年框架}账本条目记录南宋在1211年的“南宋以宋宁宗在位第17年作为诏令、赋役与外交文书的纪年框架”。其背景是新岁颁历、年号和纪年是中枢与地方行政沟通的共同前提，后续影响为为同年政令、战事和地方材料提供日期锚点，不能单独推知具体政策执行。政权、军队或制度的变化须经由道路、文书、资源和地方中介才会落实，城市、乡村、边地和流动人群不可能在同一时间承受相同结果。

本条使用《宋史》卷24—47〈本纪〉；《建炎以来系年要录》本年条；《宋会要辑稿》相关门类，并标示“纪年框架可靠；该记录仅确证政权纪年连续，年内具体事项须由本年条另证”。这些材料能够钉住年序、命令、战役或局部地点，却不能直接推出人口、税收、身份和社会动机的总量。应区分诏令与实施、条约与边境生活、军报与伤亡统计，也不将官府的族群或行政称谓写成固定的社会本质。

从区域史角度看，该事件是资源、信息与权力重新连接的节点。不同地方会因市场、交通、宗教组织、家庭策略和环境条件而产生不同反应；材料的沉默限制我们对未被记录者所能作出的判断。

\eventanchor{SS-1212-REGNAL-YEAR}{1212年·南宋以宋宁宗在位第18年作为诏令、赋役与外交文书的纪年框架}账本条目记录南宋在1212年的“南宋以宋宁宗在位第18年作为诏令、赋役与外交文书的纪年框架”。其背景是新岁颁历、年号和纪年是中枢与地方行政沟通的共同前提，后续影响为为同年政令、战事和地方材料提供日期锚点，不能单独推知具体政策执行。政权、军队或制度的变化须经由道路、文书、资源和地方中介才会落实，城市、乡村、边地和流动人群不可能在同一时间承受相同结果。

本条使用《宋史》卷24—47〈本纪〉；《建炎以来系年要录》本年条；《宋会要辑稿》相关门类，并标示“纪年框架可靠；该记录仅确证政权纪年连续，年内具体事项须由本年条另证”。这些材料能够钉住年序、命令、战役或局部地点，却不能直接推出人口、税收、身份和社会动机的总量。应区分诏令与实施、条约与边境生活、军报与伤亡统计，也不将官府的族群或行政称谓写成固定的社会本质。

从区域史角度看，该事件是资源、信息与权力重新连接的节点。不同地方会因市场、交通、宗教组织、家庭策略和环境条件而产生不同反应；材料的沉默限制我们对未被记录者所能作出的判断。

\subsection{制度、区域与材料边界}

\eventanchor{SS-1213-REGNAL-YEAR}{1213年·南宋以宋宁宗在位第19年作为诏令、赋役与外交文书的纪年框架}账本条目记录南宋在1213年的“南宋以宋宁宗在位第19年作为诏令、赋役与外交文书的纪年框架”。其背景是新岁颁历、年号和纪年是中枢与地方行政沟通的共同前提，后续影响为为同年政令、战事和地方材料提供日期锚点，不能单独推知具体政策执行。政权、军队或制度的变化须经由道路、文书、资源和地方中介才会落实，城市、乡村、边地和流动人群不可能在同一时间承受相同结果。

本条使用《宋史》卷24—47〈本纪〉；《建炎以来系年要录》本年条；《宋会要辑稿》相关门类，并标示“纪年框架可靠；该记录仅确证政权纪年连续，年内具体事项须由本年条另证”。这些材料能够钉住年序、命令、战役或局部地点，却不能直接推出人口、税收、身份和社会动机的总量。应区分诏令与实施、条约与边境生活、军报与伤亡统计，也不将官府的族群或行政称谓写成固定的社会本质。

从区域史角度看，该事件是资源、信息与权力重新连接的节点。不同地方会因市场、交通、宗教组织、家庭策略和环境条件而产生不同反应；材料的沉默限制我们对未被记录者所能作出的判断。

\eventanchor{SS-1214-REGNAL-YEAR}{1214年·南宋以宋宁宗在位第20年作为诏令、赋役与外交文书的纪年框架}账本条目记录南宋在1214年的“南宋以宋宁宗在位第20年作为诏令、赋役与外交文书的纪年框架”。其背景是新岁颁历、年号和纪年是中枢与地方行政沟通的共同前提，后续影响为为同年政令、战事和地方材料提供日期锚点，不能单独推知具体政策执行。政权、军队或制度的变化须经由道路、文书、资源和地方中介才会落实，城市、乡村、边地和流动人群不可能在同一时间承受相同结果。

本条使用《宋史》卷24—47〈本纪〉；《建炎以来系年要录》本年条；《宋会要辑稿》相关门类，并标示“纪年框架可靠；该记录仅确证政权纪年连续，年内具体事项须由本年条另证”。这些材料能够钉住年序、命令、战役或局部地点，却不能直接推出人口、税收、身份和社会动机的总量。应区分诏令与实施、条约与边境生活、军报与伤亡统计，也不将官府的族群或行政称谓写成固定的社会本质。

从区域史角度看，该事件是资源、信息与权力重新连接的节点。不同地方会因市场、交通、宗教组织、家庭策略和环境条件而产生不同反应；材料的沉默限制我们对未被记录者所能作出的判断。

\eventanchor{SS-1215-REGNAL-YEAR}{1215年·南宋以宋宁宗在位第21年作为诏令、赋役与外交文书的纪年框架}账本条目记录南宋在1215年的“南宋以宋宁宗在位第21年作为诏令、赋役与外交文书的纪年框架”。其背景是新岁颁历、年号和纪年是中枢与地方行政沟通的共同前提，后续影响为为同年政令、战事和地方材料提供日期锚点，不能单独推知具体政策执行。政权、军队或制度的变化须经由道路、文书、资源和地方中介才会落实，城市、乡村、边地和流动人群不可能在同一时间承受相同结果。

本条使用《宋史》卷24—47〈本纪〉；《建炎以来系年要录》本年条；《宋会要辑稿》相关门类，并标示“纪年框架可靠；该记录仅确证政权纪年连续，年内具体事项须由本年条另证”。这些材料能够钉住年序、命令、战役或局部地点，却不能直接推出人口、税收、身份和社会动机的总量。应区分诏令与实施、条约与边境生活、军报与伤亡统计，也不将官府的族群或行政称谓写成固定的社会本质。

从区域史角度看，该事件是资源、信息与权力重新连接的节点。不同地方会因市场、交通、宗教组织、家庭策略和环境条件而产生不同反应；材料的沉默限制我们对未被记录者所能作出的判断。

\eventanchor{SS-1216-REGNAL-YEAR}{1216年·南宋以宋宁宗在位第22年作为诏令、赋役与外交文书的纪年框架}账本条目记录南宋在1216年的“南宋以宋宁宗在位第22年作为诏令、赋役与外交文书的纪年框架”。其背景是新岁颁历、年号和纪年是中枢与地方行政沟通的共同前提，后续影响为为同年政令、战事和地方材料提供日期锚点，不能单独推知具体政策执行。政权、军队或制度的变化须经由道路、文书、资源和地方中介才会落实，城市、乡村、边地和流动人群不可能在同一时间承受相同结果。

本条使用《宋史》卷24—47〈本纪〉；《建炎以来系年要录》本年条；《宋会要辑稿》相关门类，并标示“纪年框架可靠；该记录仅确证政权纪年连续，年内具体事项须由本年条另证”。这些材料能够钉住年序、命令、战役或局部地点，却不能直接推出人口、税收、身份和社会动机的总量。应区分诏令与实施、条约与边境生活、军报与伤亡统计，也不将官府的族群或行政称谓写成固定的社会本质。

从区域史角度看，该事件是资源、信息与权力重新连接的节点。不同地方会因市场、交通、宗教组织、家庭策略和环境条件而产生不同反应；材料的沉默限制我们对未被记录者所能作出的判断。

\eventanchor{SS-1217-EVENT-135}{1217年·金军南攻四川、两淮等方向}账本条目记录南宋与金在1217年的“金军南攻四川、两淮等方向”。其背景是前期边防、财政和统治联盟的压力构成直接背景，具体条件须按本条材料分辨，后续影响为改变后续资源配置与政治选择；实际执行范围和社会影响仍须分项核验。政权、军队或制度的变化须经由道路、文书、资源和地方中介才会落实，城市、乡村、边地和流动人群不可能在同一时间承受相同结果。

本条使用《宋史》卷24—47〈本纪〉；《建炎以来系年要录》本年条；《宋会要辑稿》相关门类，并标示“编年主干可靠；细节、规模与动机须与对方材料、地方文书或考古材料互校”。这些材料能够钉住年序、命令、战役或局部地点，却不能直接推出人口、税收、身份和社会动机的总量。应区分诏令与实施、条约与边境生活、军报与伤亡统计，也不将官府的族群或行政称谓写成固定的社会本质。

从区域史角度看，该事件是资源、信息与权力重新连接的节点。不同地方会因市场、交通、宗教组织、家庭策略和环境条件而产生不同反应；材料的沉默限制我们对未被记录者所能作出的判断。

\eventanchor{SS-1217-REGNAL-YEAR}{1217年·南宋以宋宁宗在位第23年作为诏令、赋役与外交文书的纪年框架}账本条目记录南宋在1217年的“南宋以宋宁宗在位第23年作为诏令、赋役与外交文书的纪年框架”。其背景是新岁颁历、年号和纪年是中枢与地方行政沟通的共同前提，后续影响为为同年政令、战事和地方材料提供日期锚点，不能单独推知具体政策执行。政权、军队或制度的变化须经由道路、文书、资源和地方中介才会落实，城市、乡村、边地和流动人群不可能在同一时间承受相同结果。

本条使用《宋史》卷24—47〈本纪〉；《建炎以来系年要录》本年条；《宋会要辑稿》相关门类，并标示“纪年框架可靠；该记录仅确证政权纪年连续，年内具体事项须由本年条另证”。这些材料能够钉住年序、命令、战役或局部地点，却不能直接推出人口、税收、身份和社会动机的总量。应区分诏令与实施、条约与边境生活、军报与伤亡统计，也不将官府的族群或行政称谓写成固定的社会本质。

从区域史角度看，该事件是资源、信息与权力重新连接的节点。不同地方会因市场、交通、宗教组织、家庭策略和环境条件而产生不同反应；材料的沉默限制我们对未被记录者所能作出的判断。

\eventanchor{SS-1218-REGNAL-YEAR}{1218年·南宋以宋宁宗在位第24年作为诏令、赋役与外交文书的纪年框架}账本条目记录南宋在1218年的“南宋以宋宁宗在位第24年作为诏令、赋役与外交文书的纪年框架”。其背景是新岁颁历、年号和纪年是中枢与地方行政沟通的共同前提，后续影响为为同年政令、战事和地方材料提供日期锚点，不能单独推知具体政策执行。政权、军队或制度的变化须经由道路、文书、资源和地方中介才会落实，城市、乡村、边地和流动人群不可能在同一时间承受相同结果。

本条使用《宋史》卷24—47〈本纪〉；《建炎以来系年要录》本年条；《宋会要辑稿》相关门类，并标示“纪年框架可靠；该记录仅确证政权纪年连续，年内具体事项须由本年条另证”。这些材料能够钉住年序、命令、战役或局部地点，却不能直接推出人口、税收、身份和社会动机的总量。应区分诏令与实施、条约与边境生活、军报与伤亡统计，也不将官府的族群或行政称谓写成固定的社会本质。

从区域史角度看，该事件是资源、信息与权力重新连接的节点。不同地方会因市场、交通、宗教组织、家庭策略和环境条件而产生不同反应；材料的沉默限制我们对未被记录者所能作出的判断。

\eventanchor{SS-1219-REGNAL-YEAR}{1219年·南宋以宋宁宗在位第25年作为诏令、赋役与外交文书的纪年框架}账本条目记录南宋在1219年的“南宋以宋宁宗在位第25年作为诏令、赋役与外交文书的纪年框架”。其背景是新岁颁历、年号和纪年是中枢与地方行政沟通的共同前提，后续影响为为同年政令、战事和地方材料提供日期锚点，不能单独推知具体政策执行。政权、军队或制度的变化须经由道路、文书、资源和地方中介才会落实，城市、乡村、边地和流动人群不可能在同一时间承受相同结果。

本条使用《宋史》卷24—47〈本纪〉；《建炎以来系年要录》本年条；《宋会要辑稿》相关门类，并标示“纪年框架可靠；该记录仅确证政权纪年连续，年内具体事项须由本年条另证”。这些材料能够钉住年序、命令、战役或局部地点，却不能直接推出人口、税收、身份和社会动机的总量。应区分诏令与实施、条约与边境生活、军报与伤亡统计，也不将官府的族群或行政称谓写成固定的社会本质。

从区域史角度看，该事件是资源、信息与权力重新连接的节点。不同地方会因市场、交通、宗教组织、家庭策略和环境条件而产生不同反应；材料的沉默限制我们对未被记录者所能作出的判断。

\eventanchor{SS-1220-REGNAL-YEAR}{1220年·南宋以宋宁宗在位第26年作为诏令、赋役与外交文书的纪年框架}账本条目记录南宋在1220年的“南宋以宋宁宗在位第26年作为诏令、赋役与外交文书的纪年框架”。其背景是新岁颁历、年号和纪年是中枢与地方行政沟通的共同前提，后续影响为为同年政令、战事和地方材料提供日期锚点，不能单独推知具体政策执行。政权、军队或制度的变化须经由道路、文书、资源和地方中介才会落实，城市、乡村、边地和流动人群不可能在同一时间承受相同结果。

本条使用《宋史》卷24—47〈本纪〉；《建炎以来系年要录》本年条；《宋会要辑稿》相关门类，并标示“纪年框架可靠；该记录仅确证政权纪年连续，年内具体事项须由本年条另证”。这些材料能够钉住年序、命令、战役或局部地点，却不能直接推出人口、税收、身份和社会动机的总量。应区分诏令与实施、条约与边境生活、军报与伤亡统计，也不将官府的族群或行政称谓写成固定的社会本质。

从区域史角度看，该事件是资源、信息与权力重新连接的节点。不同地方会因市场、交通、宗教组织、家庭策略和环境条件而产生不同反应；材料的沉默限制我们对未被记录者所能作出的判断。

\eventanchor{SS-1221-REGNAL-YEAR}{1221年·南宋以宋宁宗在位第27年作为诏令、赋役与外交文书的纪年框架}账本条目记录南宋在1221年的“南宋以宋宁宗在位第27年作为诏令、赋役与外交文书的纪年框架”。其背景是新岁颁历、年号和纪年是中枢与地方行政沟通的共同前提，后续影响为为同年政令、战事和地方材料提供日期锚点，不能单独推知具体政策执行。政权、军队或制度的变化须经由道路、文书、资源和地方中介才会落实，城市、乡村、边地和流动人群不可能在同一时间承受相同结果。

本条使用《宋史》卷24—47〈本纪〉；《建炎以来系年要录》本年条；《宋会要辑稿》相关门类，并标示“纪年框架可靠；该记录仅确证政权纪年连续，年内具体事项须由本年条另证”。这些材料能够钉住年序、命令、战役或局部地点，却不能直接推出人口、税收、身份和社会动机的总量。应区分诏令与实施、条约与边境生活、军报与伤亡统计，也不将官府的族群或行政称谓写成固定的社会本质。

从区域史角度看，该事件是资源、信息与权力重新连接的节点。不同地方会因市场、交通、宗教组织、家庭策略和环境条件而产生不同反应；材料的沉默限制我们对未被记录者所能作出的判断。

\eventanchor{SS-1222-REGNAL-YEAR}{1222年·南宋以宋宁宗在位第28年作为诏令、赋役与外交文书的纪年框架}账本条目记录南宋在1222年的“南宋以宋宁宗在位第28年作为诏令、赋役与外交文书的纪年框架”。其背景是新岁颁历、年号和纪年是中枢与地方行政沟通的共同前提，后续影响为为同年政令、战事和地方材料提供日期锚点，不能单独推知具体政策执行。政权、军队或制度的变化须经由道路、文书、资源和地方中介才会落实，城市、乡村、边地和流动人群不可能在同一时间承受相同结果。

本条使用《宋史》卷24—47〈本纪〉；《建炎以来系年要录》本年条；《宋会要辑稿》相关门类，并标示“纪年框架可靠；该记录仅确证政权纪年连续，年内具体事项须由本年条另证”。这些材料能够钉住年序、命令、战役或局部地点，却不能直接推出人口、税收、身份和社会动机的总量。应区分诏令与实施、条约与边境生活、军报与伤亡统计，也不将官府的族群或行政称谓写成固定的社会本质。

从区域史角度看，该事件是资源、信息与权力重新连接的节点。不同地方会因市场、交通、宗教组织、家庭策略和环境条件而产生不同反应；材料的沉默限制我们对未被记录者所能作出的判断。

\eventanchor{SS-1223-REGNAL-YEAR}{1223年·南宋以宋宁宗在位第29年作为诏令、赋役与外交文书的纪年框架}账本条目记录南宋在1223年的“南宋以宋宁宗在位第29年作为诏令、赋役与外交文书的纪年框架”。其背景是新岁颁历、年号和纪年是中枢与地方行政沟通的共同前提，后续影响为为同年政令、战事和地方材料提供日期锚点，不能单独推知具体政策执行。政权、军队或制度的变化须经由道路、文书、资源和地方中介才会落实，城市、乡村、边地和流动人群不可能在同一时间承受相同结果。

本条使用《宋史》卷24—47〈本纪〉；《建炎以来系年要录》本年条；《宋会要辑稿》相关门类，并标示“纪年框架可靠；该记录仅确证政权纪年连续，年内具体事项须由本年条另证”。这些材料能够钉住年序、命令、战役或局部地点，却不能直接推出人口、税收、身份和社会动机的总量。应区分诏令与实施、条约与边境生活、军报与伤亡统计，也不将官府的族群或行政称谓写成固定的社会本质。

从区域史角度看，该事件是资源、信息与权力重新连接的节点。不同地方会因市场、交通、宗教组织、家庭策略和环境条件而产生不同反应；材料的沉默限制我们对未被记录者所能作出的判断。

\eventanchor{SS-1224-EVENT-136}{1224年·宁宗去世，理宗即位，史弥远继续主导中枢}账本条目记录南宋在1224年的“宁宗去世，理宗即位，史弥远继续主导中枢”。其背景是前期边防、财政和统治联盟的压力构成直接背景，具体条件须按本条材料分辨，后续影响为改变后续资源配置与政治选择；实际执行范围和社会影响仍须分项核验。政权、军队或制度的变化须经由道路、文书、资源和地方中介才会落实，城市、乡村、边地和流动人群不可能在同一时间承受相同结果。

本条使用《宋史》卷24—47〈本纪〉；《建炎以来系年要录》本年条；《宋会要辑稿》相关门类，并标示“编年主干可靠；细节、规模与动机须与对方材料、地方文书或考古材料互校”。这些材料能够钉住年序、命令、战役或局部地点，却不能直接推出人口、税收、身份和社会动机的总量。应区分诏令与实施、条约与边境生活、军报与伤亡统计，也不将官府的族群或行政称谓写成固定的社会本质。

从区域史角度看，该事件是资源、信息与权力重新连接的节点。不同地方会因市场、交通、宗教组织、家庭策略和环境条件而产生不同反应；材料的沉默限制我们对未被记录者所能作出的判断。

\eventanchor{SS-1224-REGNAL-YEAR}{1224年·南宋以宋宁宗在位第30年作为诏令、赋役与外交文书的纪年框架}账本条目记录南宋在1224年的“南宋以宋宁宗在位第30年作为诏令、赋役与外交文书的纪年框架”。其背景是新岁颁历、年号和纪年是中枢与地方行政沟通的共同前提，后续影响为为同年政令、战事和地方材料提供日期锚点，不能单独推知具体政策执行。政权、军队或制度的变化须经由道路、文书、资源和地方中介才会落实，城市、乡村、边地和流动人群不可能在同一时间承受相同结果。

本条使用《宋史》卷24—47〈本纪〉；《建炎以来系年要录》本年条；《宋会要辑稿》相关门类，并标示“纪年框架可靠；该记录仅确证政权纪年连续，年内具体事项须由本年条另证”。这些材料能够钉住年序、命令、战役或局部地点，却不能直接推出人口、税收、身份和社会动机的总量。应区分诏令与实施、条约与边境生活、军报与伤亡统计，也不将官府的族群或行政称谓写成固定的社会本质。

从区域史角度看，该事件是资源、信息与权力重新连接的节点。不同地方会因市场、交通、宗教组织、家庭策略和环境条件而产生不同反应；材料的沉默限制我们对未被记录者所能作出的判断。

\eventanchor{SS-1225-REGNAL-YEAR}{1225年·南宋以宋理宗在位第1年作为诏令、赋役与外交文书的纪年框架}账本条目记录南宋在1225年的“南宋以宋理宗在位第1年作为诏令、赋役与外交文书的纪年框架”。其背景是新岁颁历、年号和纪年是中枢与地方行政沟通的共同前提，后续影响为为同年政令、战事和地方材料提供日期锚点，不能单独推知具体政策执行。政权、军队或制度的变化须经由道路、文书、资源和地方中介才会落实，城市、乡村、边地和流动人群不可能在同一时间承受相同结果。

本条使用《宋史》卷24—47〈本纪〉；《建炎以来系年要录》本年条；《宋会要辑稿》相关门类，并标示“纪年框架可靠；该记录仅确证政权纪年连续，年内具体事项须由本年条另证”。这些材料能够钉住年序、命令、战役或局部地点，却不能直接推出人口、税收、身份和社会动机的总量。应区分诏令与实施、条约与边境生活、军报与伤亡统计，也不将官府的族群或行政称谓写成固定的社会本质。

从区域史角度看，该事件是资源、信息与权力重新连接的节点。不同地方会因市场、交通、宗教组织、家庭策略和环境条件而产生不同反应；材料的沉默限制我们对未被记录者所能作出的判断。

\eventanchor{SS-1226-REGNAL-YEAR}{1226年·南宋以宋理宗在位第2年作为诏令、赋役与外交文书的纪年框架}账本条目记录南宋在1226年的“南宋以宋理宗在位第2年作为诏令、赋役与外交文书的纪年框架”。其背景是新岁颁历、年号和纪年是中枢与地方行政沟通的共同前提，后续影响为为同年政令、战事和地方材料提供日期锚点，不能单独推知具体政策执行。政权、军队或制度的变化须经由道路、文书、资源和地方中介才会落实，城市、乡村、边地和流动人群不可能在同一时间承受相同结果。

本条使用《宋史》卷24—47〈本纪〉；《建炎以来系年要录》本年条；《宋会要辑稿》相关门类，并标示“纪年框架可靠；该记录仅确证政权纪年连续，年内具体事项须由本年条另证”。这些材料能够钉住年序、命令、战役或局部地点，却不能直接推出人口、税收、身份和社会动机的总量。应区分诏令与实施、条约与边境生活、军报与伤亡统计，也不将官府的族群或行政称谓写成固定的社会本质。

从区域史角度看，该事件是资源、信息与权力重新连接的节点。不同地方会因市场、交通、宗教组织、家庭策略和环境条件而产生不同反应；材料的沉默限制我们对未被记录者所能作出的判断。

\eventanchor{SS-1227-REGNAL-YEAR}{1227年·南宋以宋理宗在位第3年作为诏令、赋役与外交文书的纪年框架}账本条目记录南宋在1227年的“南宋以宋理宗在位第3年作为诏令、赋役与外交文书的纪年框架”。其背景是新岁颁历、年号和纪年是中枢与地方行政沟通的共同前提，后续影响为为同年政令、战事和地方材料提供日期锚点，不能单独推知具体政策执行。政权、军队或制度的变化须经由道路、文书、资源和地方中介才会落实，城市、乡村、边地和流动人群不可能在同一时间承受相同结果。

本条使用《宋史》卷24—47〈本纪〉；《建炎以来系年要录》本年条；《宋会要辑稿》相关门类，并标示“纪年框架可靠；该记录仅确证政权纪年连续，年内具体事项须由本年条另证”。这些材料能够钉住年序、命令、战役或局部地点，却不能直接推出人口、税收、身份和社会动机的总量。应区分诏令与实施、条约与边境生活、军报与伤亡统计，也不将官府的族群或行政称谓写成固定的社会本质。

从区域史角度看，该事件是资源、信息与权力重新连接的节点。不同地方会因市场、交通、宗教组织、家庭策略和环境条件而产生不同反应；材料的沉默限制我们对未被记录者所能作出的判断。

\eventanchor{SS-1228-REGNAL-YEAR}{1228年·南宋以宋理宗在位第4年作为诏令、赋役与外交文书的纪年框架}账本条目记录南宋在1228年的“南宋以宋理宗在位第4年作为诏令、赋役与外交文书的纪年框架”。其背景是新岁颁历、年号和纪年是中枢与地方行政沟通的共同前提，后续影响为为同年政令、战事和地方材料提供日期锚点，不能单独推知具体政策执行。政权、军队或制度的变化须经由道路、文书、资源和地方中介才会落实，城市、乡村、边地和流动人群不可能在同一时间承受相同结果。

本条使用《宋史》卷24—47〈本纪〉；《建炎以来系年要录》本年条；《宋会要辑稿》相关门类，并标示“纪年框架可靠；该记录仅确证政权纪年连续，年内具体事项须由本年条另证”。这些材料能够钉住年序、命令、战役或局部地点，却不能直接推出人口、税收、身份和社会动机的总量。应区分诏令与实施、条约与边境生活、军报与伤亡统计，也不将官府的族群或行政称谓写成固定的社会本质。

从区域史角度看，该事件是资源、信息与权力重新连接的节点。不同地方会因市场、交通、宗教组织、家庭策略和环境条件而产生不同反应；材料的沉默限制我们对未被记录者所能作出的判断。

\eventanchor{SS-1229-REGNAL-YEAR}{1229年·南宋以宋理宗在位第5年作为诏令、赋役与外交文书的纪年框架}账本条目记录南宋在1229年的“南宋以宋理宗在位第5年作为诏令、赋役与外交文书的纪年框架”。其背景是新岁颁历、年号和纪年是中枢与地方行政沟通的共同前提，后续影响为为同年政令、战事和地方材料提供日期锚点，不能单独推知具体政策执行。政权、军队或制度的变化须经由道路、文书、资源和地方中介才会落实，城市、乡村、边地和流动人群不可能在同一时间承受相同结果。

本条使用《宋史》卷24—47〈本纪〉；《建炎以来系年要录》本年条；《宋会要辑稿》相关门类，并标示“纪年框架可靠；该记录仅确证政权纪年连续，年内具体事项须由本年条另证”。这些材料能够钉住年序、命令、战役或局部地点，却不能直接推出人口、税收、身份和社会动机的总量。应区分诏令与实施、条约与边境生活、军报与伤亡统计，也不将官府的族群或行政称谓写成固定的社会本质。

从区域史角度看，该事件是资源、信息与权力重新连接的节点。不同地方会因市场、交通、宗教组织、家庭策略和环境条件而产生不同反应；材料的沉默限制我们对未被记录者所能作出的判断。

\eventanchor{SS-1230-REGNAL-YEAR}{1230年·南宋以宋理宗在位第6年作为诏令、赋役与外交文书的纪年框架}账本条目记录南宋在1230年的“南宋以宋理宗在位第6年作为诏令、赋役与外交文书的纪年框架”。其背景是新岁颁历、年号和纪年是中枢与地方行政沟通的共同前提，后续影响为为同年政令、战事和地方材料提供日期锚点，不能单独推知具体政策执行。政权、军队或制度的变化须经由道路、文书、资源和地方中介才会落实，城市、乡村、边地和流动人群不可能在同一时间承受相同结果。

本条使用《宋史》卷24—47〈本纪〉；《建炎以来系年要录》本年条；《宋会要辑稿》相关门类，并标示“纪年框架可靠；该记录仅确证政权纪年连续，年内具体事项须由本年条另证”。这些材料能够钉住年序、命令、战役或局部地点，却不能直接推出人口、税收、身份和社会动机的总量。应区分诏令与实施、条约与边境生活、军报与伤亡统计，也不将官府的族群或行政称谓写成固定的社会本质。

从区域史角度看，该事件是资源、信息与权力重新连接的节点。不同地方会因市场、交通、宗教组织、家庭策略和环境条件而产生不同反应；材料的沉默限制我们对未被记录者所能作出的判断。

\eventanchor{SS-1231-REGNAL-YEAR}{1231年·南宋以宋理宗在位第7年作为诏令、赋役与外交文书的纪年框架}账本条目记录南宋在1231年的“南宋以宋理宗在位第7年作为诏令、赋役与外交文书的纪年框架”。其背景是新岁颁历、年号和纪年是中枢与地方行政沟通的共同前提，后续影响为为同年政令、战事和地方材料提供日期锚点，不能单独推知具体政策执行。政权、军队或制度的变化须经由道路、文书、资源和地方中介才会落实，城市、乡村、边地和流动人群不可能在同一时间承受相同结果。

本条使用《宋史》卷24—47〈本纪〉；《建炎以来系年要录》本年条；《宋会要辑稿》相关门类，并标示“纪年框架可靠；该记录仅确证政权纪年连续，年内具体事项须由本年条另证”。这些材料能够钉住年序、命令、战役或局部地点，却不能直接推出人口、税收、身份和社会动机的总量。应区分诏令与实施、条约与边境生活、军报与伤亡统计，也不将官府的族群或行政称谓写成固定的社会本质。

从区域史角度看，该事件是资源、信息与权力重新连接的节点。不同地方会因市场、交通、宗教组织、家庭策略和环境条件而产生不同反应；材料的沉默限制我们对未被记录者所能作出的判断。

\eventanchor{SS-1232-REGNAL-YEAR}{1232年·南宋以宋理宗在位第8年作为诏令、赋役与外交文书的纪年框架}账本条目记录南宋在1232年的“南宋以宋理宗在位第8年作为诏令、赋役与外交文书的纪年框架”。其背景是新岁颁历、年号和纪年是中枢与地方行政沟通的共同前提，后续影响为为同年政令、战事和地方材料提供日期锚点，不能单独推知具体政策执行。政权、军队或制度的变化须经由道路、文书、资源和地方中介才会落实，城市、乡村、边地和流动人群不可能在同一时间承受相同结果。

本条使用《宋史》卷24—47〈本纪〉；《建炎以来系年要录》本年条；《宋会要辑稿》相关门类，并标示“纪年框架可靠；该记录仅确证政权纪年连续，年内具体事项须由本年条另证”。这些材料能够钉住年序、命令、战役或局部地点，却不能直接推出人口、税收、身份和社会动机的总量。应区分诏令与实施、条约与边境生活、军报与伤亡统计，也不将官府的族群或行政称谓写成固定的社会本质。

从区域史角度看，该事件是资源、信息与权力重新连接的节点。不同地方会因市场、交通、宗教组织、家庭策略和环境条件而产生不同反应；材料的沉默限制我们对未被记录者所能作出的判断。

\eventanchor{SS-1233-EVENT-137}{1233年·南宋与蒙古在攻金问题上形成短暂军事协作}账本条目记录南宋与蒙古在1233年的“南宋与蒙古在攻金问题上形成短暂军事协作”。其背景是前期边防、财政和统治联盟的压力构成直接背景，具体条件须按本条材料分辨，后续影响为改变后续资源配置与政治选择；实际执行范围和社会影响仍须分项核验。政权、军队或制度的变化须经由道路、文书、资源和地方中介才会落实，城市、乡村、边地和流动人群不可能在同一时间承受相同结果。

本条使用《宋史》卷24—47〈本纪〉；《建炎以来系年要录》本年条；《宋会要辑稿》相关门类，并标示“编年主干可靠；细节、规模与动机须与对方材料、地方文书或考古材料互校”。这些材料能够钉住年序、命令、战役或局部地点，却不能直接推出人口、税收、身份和社会动机的总量。应区分诏令与实施、条约与边境生活、军报与伤亡统计，也不将官府的族群或行政称谓写成固定的社会本质。

从区域史角度看，该事件是资源、信息与权力重新连接的节点。不同地方会因市场、交通、宗教组织、家庭策略和环境条件而产生不同反应；材料的沉默限制我们对未被记录者所能作出的判断。

\eventanchor{SS-1233-REGNAL-YEAR}{1233年·南宋以宋理宗在位第9年作为诏令、赋役与外交文书的纪年框架}账本条目记录南宋在1233年的“南宋以宋理宗在位第9年作为诏令、赋役与外交文书的纪年框架”。其背景是新岁颁历、年号和纪年是中枢与地方行政沟通的共同前提，后续影响为为同年政令、战事和地方材料提供日期锚点，不能单独推知具体政策执行。政权、军队或制度的变化须经由道路、文书、资源和地方中介才会落实，城市、乡村、边地和流动人群不可能在同一时间承受相同结果。

本条使用《宋史》卷24—47〈本纪〉；《建炎以来系年要录》本年条；《宋会要辑稿》相关门类，并标示“纪年框架可靠；该记录仅确证政权纪年连续，年内具体事项须由本年条另证”。这些材料能够钉住年序、命令、战役或局部地点，却不能直接推出人口、税收、身份和社会动机的总量。应区分诏令与实施、条约与边境生活、军报与伤亡统计，也不将官府的族群或行政称谓写成固定的社会本质。

从区域史角度看，该事件是资源、信息与权力重新连接的节点。不同地方会因市场、交通、宗教组织、家庭策略和环境条件而产生不同反应；材料的沉默限制我们对未被记录者所能作出的判断。

\eventanchor{SS-1234-EVENT-138}{1234年·蔡州陷落后，南宋试图进入河南故地}账本条目记录南宋金蒙古在1234年的“蔡州陷落后，南宋试图进入河南故地”。其背景是前期边防、财政和统治联盟的压力构成直接背景，具体条件须按本条材料分辨，后续影响为改变后续资源配置与政治选择；实际执行范围和社会影响仍须分项核验。政权、军队或制度的变化须经由道路、文书、资源和地方中介才会落实，城市、乡村、边地和流动人群不可能在同一时间承受相同结果。

本条使用《宋史》卷24—47〈本纪〉；《建炎以来系年要录》本年条；《宋会要辑稿》相关门类，并标示“编年主干可靠；细节、规模与动机须与对方材料、地方文书或考古材料互校”。这些材料能够钉住年序、命令、战役或局部地点，却不能直接推出人口、税收、身份和社会动机的总量。应区分诏令与实施、条约与边境生活、军报与伤亡统计，也不将官府的族群或行政称谓写成固定的社会本质。

从区域史角度看，该事件是资源、信息与权力重新连接的节点。不同地方会因市场、交通、宗教组织、家庭策略和环境条件而产生不同反应；材料的沉默限制我们对未被记录者所能作出的判断。
